% Requires: amsmath, amssymb, listings, booktabs, array (assumed in parent document)
\section{QRDQN en Stable Baselines3 para el TFG}

\subsection{Resumen ejecutivo}

La implementación de \textbf{QR-DQN} en el ecosistema de Stable Baselines vive en
\textbf{SB3-Contrib}, no en el núcleo de Stable Baselines3, y convierte a DQN en un método
\textbf{distribucional}: en vez de aprender un único escalar \(Q(s,a)\), aprende una
\textbf{distribución aproximada de retornos} por acción usando \(N\) cuantiles. En la práctica de
SB3, esa distribución se representa como un tensor de forma
\((\text{batch},\, N,\, |\mathcal{A}|)\), se entrena con \textbf{quantile Huber loss}, y la acción
codiciosa se elige por la \textbf{media de los cuantiles}, es decir, por el retorno esperado
implícito en la distribución. Por tanto, QRDQN aprende más información que DQN, pero \textbf{no es
risk-sensitive por defecto}: sigue actuando según la esperanza matemática salvo que se modifique
explícitamente el criterio de decisión.

Desde la perspectiva teórica, QR-DQN nace para cerrar la brecha entre la teoría distribucional de
Marc~G.~Bellemare, Will~Dabney, y Rémi~Munos en C51 y una implementación entrenable por descenso de
gradiente. C51 fija los soportes \(z_i\) y aprende probabilidades; QR-DQN ``transpone'' esa idea:
fija probabilidades uniformes y aprende las \textbf{localizaciones} de los cuantiles. Esto evita la
proyección categórica de C51 y elimina la necesidad de fijar manualmente cotas como \(V_{\min}\) y
\(V_{\max}\).

En el repositorio \texttt{yeaight7/TFG\_CYBER\_AI}, QRDQN es el algoritmo principal para un
problema discreto binario de ciberdefensa sobre \textbf{CICIDS2017}, con observaciones de
\textbf{152 dimensiones} y acciones \textbf{0=PERMIT}, \textbf{1=BLOCK}. La arquitectura elegida es
\texttt{MlpPolicy} con \texttt{net\_arch=[512, 256]}, \texttt{learning\_rate=1e-4},
\texttt{gamma=0.99}, \texttt{tau=1.0}, \texttt{train\_freq} de 50 o 100 y
\texttt{gradient\_steps} de 10 o 20 según preset. El mejor artefacto consolidado del repositorio es
el run \texttt{C03\_qrdqn\_cicids2017\_canonical\_full\_random\_20260223\_232439}, con accuracy
\texttt{0.99859} y F1 de ataque \texttt{0.99876} en random split; sin embargo, la validación dura
por día/CSV cae de forma importante, lo que demuestra que la generalización real es bastante más
difícil que el escenario aleatorio.

La conclusión práctica para el TFG es clara: QRDQN está bien alineado con el problema porque el
espacio de acciones es discreto y porque el aprendizaje distribucional puede capturar mejor colas,
multimodalidad y asimetrías del retorno inducidas por una recompensa con falsos negativos muy
penalizados. Pero, en SB3, esa ventaja depende mucho de detalles de implementación:
\textbf{número de cuantiles}, rapidez del decaimiento de \(\varepsilon\), tamaño del batch,
frecuencia de actualización de la red objetivo y diseño de la recompensa.

\subsection{Supuestos y lagunas del repositorio}

\begin{itemize}
  \item \textbf{La versión exacta de SB3 y SB3-Contrib no está fijada}, solo un mínimo:
        \texttt{stable-baselines3>=2.3} y \texttt{sb3-contrib>=2.3}. En SB3-Contrib v2.3.0 cambió
        el default de \texttt{learning\_starts} de 50\,000 a 100.

  \item \textbf{\texttt{n\_quantiles} no se fija explícitamente} en el script de entrenamiento; se
        infiere que se usa el default oficial: \textbf{200 cuantiles}.

  \item \textbf{Los parámetros de exploración tampoco se fijan explícitamente}. Se heredan los
        defaults de QRDQN en SB3-Contrib: \texttt{exploration\_fraction=0.005},
        \texttt{exploration\_initial\_eps=1.0} y \texttt{exploration\_final\_eps=0.01}. Con
        100\,000 timesteps, la caída lineal de \(\varepsilon\) termina aproximadamente en los
        \textbf{primeros 500 pasos}; con 25\,000 timesteps, en los \textbf{primeros 125 pasos}.

  \item \textbf{El script \verb|train_rl_defender.py| contiene una inconsistencia documental
        interna}: en la ayuda de \texttt{--timesteps} se habla de un default de ``500k full, 5k
        fast'', pero el código efectivo en \texttt{main()} asigna \texttt{25\_000} para
        \texttt{fast} y \texttt{100\_000} para \texttt{full}.

  \item \textbf{Los resultados históricos no coinciden necesariamente con los defaults actuales del
        código}. El reward actual por defecto usa \texttt{fp=-1.5}, mientras que el mejor run
        histórico consolidado usó \texttt{fp=-2.0}.

  \item \textbf{La validación leave-one-exact-CSV-out está implementada, pero no hay un artefacto
        agregado comprometido} en \verb|runs/validation/|.
\end{itemize}

Como hallazgo sintético del repo: el pipeline actual trabaja con CICIDS2017, un esquema canónico de
\textbf{76 features}, una máscara de missingness de otras \textbf{76}, observación total de
\textbf{152 dimensiones}, y un entorno Gymnasium con acción discreta binaria.

\subsection{Fundamentos teóricos y matemáticos}

La motivación de \textbf{distributional RL} es que el retorno futuro desde un estado-acción no es
un número fijo, sino una \textbf{variable aleatoria}. DQN clásico aprende solo su expectativa
\(Q(s,a)=\mathbb{E}[Z(s,a)]\). El enfoque distribucional intenta aproximar la ley completa de
\(Z(s,a)\), lo que permite representar colas, multimodalidad y asimetría.

En C51, la distribución se aproxima con una combinación discreta de \textbf{soportes fijos}
\(z_1,\dots,z_N\) y \textbf{probabilidades aprendidas} \(q_i\). QR-DQN invierte esa
parametrización: fija probabilidades uniformes \(1/N\) y aprende las \textbf{localizaciones}
\(\theta_i(s,a)\). Formalmente:
\[
  Z_\theta(s,a)=\frac{1}{N}\sum_{i=1}^N \delta_{\theta_i(s,a)},
\]
donde cada \(\theta_i\) representa una localización cuantil y \(\delta\) es una masa de Dirac.

La \textbf{pinball loss} para un único cuantil \(\tau\):
\[
  \rho_\tau(u)=u\bigl(\tau-\mathbf{1}_{\{u<0\}}\bigr).
\]

Como la pinball loss no es suave en cero, el paper propone la \textbf{quantile Huber loss}:
\[
  L_\kappa(u)=
  \begin{cases}
    \tfrac{1}{2}u^2, & |u|\le \kappa \\
    \kappa\bigl(|u|-\tfrac{1}{2}\kappa\bigr), & |u|>\kappa
  \end{cases}
  \qquad
  \rho^\kappa_\tau(u)=|\tau-\mathbf{1}_{\{u<0\}}|\,L_\kappa(u).
\]

El \textbf{Bellman update} distribucional en control:
\[
  \mathcal{T} Z(s,a)\overset{D}{=} r + \gamma Z(s',a^*),
  \qquad
  a^*=\arg\max_{a'}\mathbb{E}[Z(s',a')].
\]

\begin{table}[htbp]
\caption{Comparación entre DQN, C51 y QRDQN.}
\footnotesize
\begin{tabular}{p{3.5cm}p{3cm}p{3cm}p{4cm}}
\hline
\textbf{Atributo} & \textbf{DQN} & \textbf{C51} & \textbf{QRDQN} \\
\hline
Qué modela
  & Esperanza \(Q(s,a)\)
  & Distribución categórica
  & Distribución por cuantiles \\
Representación
  & 1 escalar por acción
  & Soportes fijos + probs aprendidas
  & Probs uniformes + localizaciones aprendidas \\
Salida por acción
  & 1 valor
  & \(N\) probabilidades
  & \(N\) cuantiles \\
Necesita \(V_{\min},V_{\max}\)
  & No & Sí & No \\
Pérdida
  & TD/Huber escalar
  & KL a proyección categórica
  & Quantile Huber \\
Política greedy
  & \(\arg\max Q\)
  & \(\arg\max \mathbb{E}[Z]\)
  & \(\arg\max \mathbb{E}[Z]\) \\
Ventaja principal
  & Simplicidad
  & Distribucional con coste moderado
  & Distribucional más flexible, sin soportes fijos \\
Inconveniente principal
  & Pierde información de dispersión
  & Requiere proyección y cotas
  & Cabeza de salida más grande y más coste \\
\hline
\end{tabular}
\end{table}

\subsection{Lógica algorítmica e implementación en SB3}

A nivel algorítmico, la lógica de QRDQN en SB3 es la de un \textbf{off-policy method} clásico:
recolecta transiciones en un replay buffer, espera hasta \texttt{learning\_starts}, entrena cada
\texttt{train\_freq} pasos durante \texttt{gradient\_steps} iteraciones y actualiza la red objetivo
cada \texttt{target\_update\_interval}.

% Diagrama original en Mermaid (flowchart TD):
\begin{verbatim}
obs_t
 -> selección de acción
 -> env.step
 -> guardar transición en replay buffer
 -> [num_timesteps > learning_starts?]
     no  -> obs_t
     sí  -> [se alcanzó train_freq?]
             no  -> obs_t
             sí  -> muestra minibatch
                    -> Z_next = target_net(next_obs)
                    -> a* = argmax_a media_tau Z_next
                    -> Z_target = r + gamma * (1-done) * Z_next[a*]
                    -> Z_current = online_net(obs, action)
                    -> quantile_huber_loss
                    -> Adam / backprop
                    -> [target_update_interval?]
                        sí -> polyak_update -> obs_t
                        no -> obs_t
\end{verbatim}

El detalle más importante de la implementación oficial es \textbf{cómo se elige la acción
bootstrap}. SB3 calcula los cuantiles del siguiente estado con la red objetivo, toma la
\textbf{media} sobre el eje de cuantiles, aplica \texttt{argmax} sobre acciones y luego recoge los
cuantiles de esa acción ganadora para formar el target.

\begin{lstlisting}[language=Python, caption={Pseudocódigo del entrenamiento QRDQN en SB3-Contrib}]
# Pseudocodigo fiel al entrenamiento de SB3-Contrib QRDQN
Z_next = target_net(next_obs)                  # (B, N, A)
a_star = argmax_a(mean_tau(Z_next))            # accion greedy por esperanza
Z_boot = select_action_quantiles(Z_next, a_star)  # (B, N)
Z_target = r + (1 - done) * gamma * Z_boot

Z_current = select_action_quantiles(online_net(obs), action)
loss = quantile_huber_loss(Z_current, Z_target)
\end{lstlisting}

En cuanto al \textbf{replay buffer}, QRDQN hereda la lógica genérica de
\texttt{OffPolicyAlgorithm}: \textbf{1-step targets} y replay uniforme, \textbf{sin prioritized
replay nativo}.

La \textbf{red objetivo} se actualiza con \texttt{polyak\_update}. Con \texttt{tau=1.0} (valor
usado en el repo) esto equivale a una \textbf{copia dura} cada
\texttt{target\_update\_interval}.

La \textbf{exploración} es \(\varepsilon\)-greedy con programación lineal entre
\texttt{exploration\_initial\_eps} y \texttt{exploration\_final\_eps} a lo largo de
\texttt{exploration\_fraction}.

El \textbf{optimizador} por defecto es \texttt{Adam}. QRDQN inyecta
\texttt{optimizer\_kwargs=\{"eps": 0.01 / batch\_size\}}. Con \texttt{batch\_size=2048}, el
\(\varepsilon\) implícito del Adam es $\approx 4.88\times10^{-6}$; con \texttt{batch\_size=512},
$\approx 1.95\times10^{-5}$.

\subsection{Arquitectura de red y parámetros de la API}

La política oficial \texttt{QRDQNPolicy} construye dos redes: \texttt{quantile\_net} y
\texttt{quantile\_net\_target}. Para observaciones vectoriales, el extractor por defecto es
\texttt{FlattenExtractor} y la arquitectura MLP por defecto es \texttt{[64, 64]}. La capa final
genera \texttt{action\_dim * n\_quantiles} salidas reordenadas a forma
\texttt{(-1, n\_quantiles, n\_actions)}.

% Diagrama original en Mermaid (flowchart LR):
\begin{verbatim}
Observacion x -> FlattenExtractor -> MLP hidden 1 -> MLP hidden 2
  -> Linear(A*N) -> reshape(batch, N, A) -> media sobre N -> argmax accion
\end{verbatim}

En el repo, la arquitectura concreta es \texttt{MlpPolicy} con \texttt{net\_arch=[512, 256]}. Con
\texttt{n\_quantiles=200} (default), el cabezal de salida tiene \textbf{400 neuronas}
$(2\,\text{acciones} \times 200\,\text{cuantiles})$. La forma de salida es $(B,\, 200,\, 2)$.

El recuento de parámetros de la red online para una MLP vectorial es:
\[
  (d h_1 + h_1) + (h_1 h_2 + h_2) + (h_2 A N + A N),
\]
donde \(d\) es la dimensión de observación, \(A\) el número de acciones y \(N\) el número de
cuantiles. Con \(d=152\), \(A=2\), \(N=200\), \(h_1=512\), \(h_2=256\), el total es
\textbf{312\,464 parámetros} en la red online.

\subsubsection{Parámetros del constructor \texttt{QRDQN}}

\begin{table}[htbp]
\caption{Parámetros del constructor \texttt{QRDQN} y su estado en el repositorio.}
\scriptsize
\begin{tabular}{p{2.8cm}p{1.8cm}p{3cm}p{3.5cm}p{2.5cm}}
\hline
\textbf{Parámetro} & \textbf{Default} & \textbf{Papel} & \textbf{Valores típicos / efecto}
  & \textbf{En el repo} \\
\hline
\texttt{policy}
  & ---
  & tipo de política
  & depende de observación
  & \textbf{Sí}: \texttt{"MlpPolicy"} \\
\texttt{env}
  & ---
  & entorno Gym/Gymnasium o VecEnv
  & obligatorio
  & \textbf{Sí}: \texttt{DummyVecEnv(Monitor(...))} \\
\texttt{learning\_rate}
  & \texttt{5e-5}
  & ritmo de aprendizaje
  & más alto acelera, pero arriesga inestabilidad
  & \textbf{Sí}: \texttt{1e-4} \\
\texttt{buffer\_size}
  & \texttt{1\_000\_000}
  & tamaño máximo del replay buffer
  & buffers grandes estabilizan, consumen RAM
  & \textbf{Sí}: cap en 200k \\
\texttt{learning\_starts}
  & \texttt{100}
  & warm-up antes de entrenar
  & demasiado bajo: targets ruidosos
  & \textbf{Implícito}: default 100 \\
\texttt{batch\_size}
  & \texttt{32}
  & tamaño de minibatch
  & mayor batch reduce varianza
  & \textbf{Sí}: \texttt{512} / \texttt{2048} \\
\texttt{tau}
  & \texttt{1.0}
  & soft update de target; \texttt{1.0} = hard copy
  & \texttt{<1} blando; \texttt{1} duro
  & \textbf{Sí}: \texttt{1.0} \\
\texttt{gamma}
  & \texttt{0.99}
  & descuento
  & más alto: horizonte largo
  & \textbf{Sí}: \texttt{0.99} \\
\texttt{train\_freq}
  & \texttt{4}
  & frecuencia de updates
  & controla cadencia entre datos y gradientes
  & \textbf{Sí}: \texttt{50} / \texttt{100} \\
\texttt{gradient\_steps}
  & \texttt{1}
  & nº de pasos de gradiente por ciclo
  & más alto = más compute
  & \textbf{Sí}: \texttt{10} / \texttt{20} \\
\texttt{replay\_buffer\_class}
  & \texttt{None}
  & clase de buffer
  & custom buffer si se quiere extender
  & \textbf{No} \\
\texttt{optimize\_memory\_usage}
  & \texttt{False}
  & variante memory-efficient
  & ahorra memoria, añade complejidad
  & \textbf{No} \\
\texttt{target\_update\_interval}
  & \texttt{10000}
  & cada cuántos pasos se actualiza target
  & pequeño = más reactivo; grande = más estable
  & \textbf{Sí}: \texttt{1000} / \texttt{10000} \\
\texttt{exploration\_fraction}
  & \texttt{0.005}
  & fracción donde cae \(\varepsilon\)
  & más grande = exploración prolongada
  & \textbf{Implícito} \\
\texttt{exploration\_initial\_eps}
  & \texttt{1.0}
  & \(\varepsilon\) inicial
  & normalmente alto
  & \textbf{Implícito} \\
\texttt{exploration\_final\_eps}
  & \texttt{0.01}
  & \(\varepsilon\) final
  & más alto = más exploración residual
  & \textbf{Implícito} \\
\texttt{policy\_kwargs}
  & \texttt{None}
  & kwargs internos de la política
  & donde se pasa \texttt{net\_arch}, \texttt{n\_quantiles}
  & \textbf{Sí} \\
\texttt{seed}
  & \texttt{None}
  & semilla global
  & clave para reproducibilidad local
  & \textbf{Sí}: \texttt{42} por defecto CLI \\
\texttt{device}
  & \texttt{"auto"}
  & CPU / CUDA
  & afecta velocidad y reproducibilidad
  & \textbf{Sí}: CPU o CUDA detectado \\
\hline
\end{tabular}
\end{table}

\subsubsection{Parámetros internos de \texttt{QRDQNPolicy}}

\begin{table}[htbp]
\caption{Parámetros de \texttt{QRDQNPolicy} propagados vía \texttt{policy\_kwargs}.}
\footnotesize
\begin{tabular}{p{3.5cm}p{3.5cm}p{3.5cm}p{3cm}}
\hline
\textbf{Parámetro} & \textbf{Default oficial} & \textbf{Papel} & \textbf{En el repo} \\
\hline
\texttt{n\_quantiles}
  & \texttt{200}
  & nº de cuantiles por acción; afecta resolución y coste
  & \textbf{Implícito} \\
\texttt{net\_arch}
  & \texttt{[64,64]} en MLP
  & profundidad/anchura MLP
  & \textbf{Sí}: \texttt{[512,256]} \\
\texttt{activation\_fn}
  & \texttt{ReLU}
  & no linealidad
  & \textbf{Implícito} \\
\texttt{features\_extractor\_class}
  & \texttt{FlattenExtractor}
  & extractor para observación vectorial
  & \textbf{Implícito} \\
\texttt{normalize\_images}
  & \texttt{True}
  & relevante en observaciones de imagen
  & \textbf{Implícito}, irrelevante aquí \\
\texttt{optimizer\_class}
  & \texttt{Adam}
  & optimizador
  & \textbf{Implícito} \\
\texttt{optimizer\_kwargs}
  & QRDQN inyecta \texttt{eps=0.01/batch\_size}
  & detalle fino de estabilidad
  & \textbf{Implícito} \\
\hline
\end{tabular}
\end{table}

\subsubsection{Fragmento relevante del repositorio}

\begin{lstlisting}[language=Python]
model = QRDQN(
    "MlpPolicy",
    vec_env,
    seed=seed,
    policy_kwargs=dict(net_arch=[512, 256]),
    learning_rate=lr,
    buffer_size=min(200_000, max(total_timesteps, 10_000)),
    batch_size=batch_size,
    gradient_steps=gradient_steps,
    gamma=0.99,
    tau=1.0,
    train_freq=train_freq,
    target_update_interval=target_update_interval,
    verbose=1,
    device=device,
    tensorboard_log=tb_log_dir,
)
\end{lstlisting}

\subsection{Consideraciones prácticas, evaluación y reproducibilidad}

QRDQN suele ser una buena elección frente a DQN cuando el problema tiene \textbf{acciones
discretas} y se sospecha que el retorno no está bien descrito por una media. Frente a C51, QRDQN
resulta especialmente cómodo cuando no se quieren elegir manualmente cotas de soporte.

En SB3, sin embargo, esa superioridad no es automática. La guía oficial recomienda hacer tuning
cuantitativo, evaluar en un entorno separado y validar muy bien terminaciones y wrappers.

El gran resultado positivo del repo es el rendimiento en random split:
\texttt{C03\_qrdqn\_cicids2017\_canonical\_full\_random\_20260223\_232439} reporta accuracy
\texttt{0.99859} y F1 de ataque \texttt{0.99876}. El hallazgo científicamente más valioso es
negativo: en la validación dura por día/CSV (Check~C) la accuracy baja a \texttt{0.84135} y el
recall de ataque a \texttt{0.52954}. Esto permite defender algo muy sólido: \textbf{el modelo
aprende muy bien dentro del régimen i.i.d. aproximado del split aleatorio, pero generalizar entre
días y capturas es mucho más difícil}.

Los \textbf{dos hiperparámetros críticos} a revisar en primer lugar son:
\begin{itemize}
  \item \texttt{exploration\_fraction}: al heredarse el default \texttt{0.005} puede que se explore
        demasiado poco durante casi todo el entrenamiento.
  \item \texttt{n\_quantiles}: \texttt{200} es un default razonable en Atari pero quizá no sea
        óptimo para un problema tabular/vectorial con dos acciones.
\end{itemize}

Para visualizar la salida distribucional en una defensa:

\begin{verbatim}
Accion PERMIT:
q05 = -3.2    q50 = 0.1    q95 = 1.7    media ≈ -0.2

Accion BLOCK:
q05 = -0.9    q50 = 0.8    q95 = 1.1    media ≈  0.7

Decision greedy de SB3: BLOCK
Interpretacion: BLOCK tiene mejor retorno esperado y menor dispersion.
\end{verbatim}

Utilidad mínima para extraer curvas de cuantiles programáticamente:

\begin{lstlisting}[language=Python]
import torch
import numpy as np

def get_action_quantiles(model, obs_np):
    obs_t = torch.as_tensor(obs_np[None], device=model.device, dtype=torch.float32)
    with torch.no_grad():
        z = model.quantile_net(obs_t)   # (1, N, A)
    z = z.squeeze(0).cpu().numpy()      # (N, A)
    q_mean = z.mean(axis=0)             # (A,)
    return z, q_mean
\end{lstlisting}

\subsection{Preguntas de tribunal y respuestas breves}

\begin{description}
  \item[¿Qué aprende QRDQN que DQN no aprende?]
    DQN aprende solo el retorno esperado por acción. QRDQN aprende una aproximación de la
    \textbf{distribución completa de retornos} mediante \(N\) cuantiles por acción.

  \item[¿En qué se diferencia QRDQN de C51 en una frase?]
    C51 fija los soportes y aprende probabilidades; QRDQN fija probabilidades uniformes y aprende
    las localizaciones de los cuantiles. Por eso QRDQN no necesita \(V_{\min}\) ni \(V_{\max}\).

  \item[Si QRDQN aprende distribución, ¿por qué no es ``risk-sensitive''?]
    Porque en SB3 la acción se elige con \texttt{argmax} de la \textbf{media} de cuantiles. La
    distribución se aprende y se puede analizar, pero la política por defecto sigue maximizando
    esperanza.

  \item[¿Tu implementación usa prioritized replay, Double DQN o dueling?]
    No de forma nativa. El QRDQN de SB3-Contrib usa \texttt{ReplayBuffer} estándar y muestreo
    uniforme.

  \item[¿Cuál es el hiperparámetro más ``silencioso'' de tu implementación?]
    Probablemente \texttt{exploration\_fraction}, porque no se fija en el repo y el default
    \texttt{0.005} implica que \(\varepsilon\) cae muy deprisa.

  \item[¿Por qué no basta con reportar accuracy de random split?]
    Porque el propio repositorio muestra que el rendimiento en split aleatorio es casi perfecto,
    pero cae de forma notable en el split duro por día/CSV. Eso indica que el verdadero reto es la
    \textbf{generalización fuera de distribución}.

  \item[¿Cuál sería una mejora inmediata para la memoria del TFG?]
    Fijar en texto y código los parámetros implícitos heredados de SB3-Contrib: versión exacta,
    \texttt{n\_quantiles}, exploración y \texttt{optimizer\_kwargs}.
\end{description}
